\par\Needspace{8\baselineskip}\bigskip
\pdfbookmark[3]{Darstellung, 5}{pallas-II-article5}
\begin{center}5.\end{center}
\noindent\hspace*{\parindent}Es ist angebracht, hier ein Lemma einzufügen, das wir nirgends formuliert gefunden haben und das sich oft mit Vorteil verwenden lässt.
\par\textit{Seien \(Q,R,Q',R'\) vier Punkte auf der Oberfläche einer Kugel und \(\theta\) der Winkel zwischen den Großkreisen \(QR,Q'R'\). Dabei betrachtet man die beiden Äste, die von ihrem Schnittpunkt aus in den Richtungen von \(Q\) nach \(R\) und von \(Q'\) nach \(R'\) verlaufen. Dann gilt}
\[
\begin{aligned}
&\cos QQ'\mathbin{.}\cos RR'-\cos QR'\mathbin{.}\cos RQ'\\
&\qquad=\cos\theta\mathbin{.}\sin QR\mathbin{.}\sin Q'R'.
\end{aligned}
\]
\par Man kann an die Stelle von \(\theta\) den Bogen \(SS'\) setzen, wobei \(S\) und \(S'\) die Pole der Großkreise \(QR,Q'R'\) bezeichnen und diese Pole in gleicher Weise zu den genannten \SourceLeaf{453}Großkreisen liegen. Das heißt, die Bewegung von \(Q\) nach \(R\) um \(S\) hat denselben Drehsinn wie die Bewegung von \(Q'\) nach \(R'\) um \(S'\).
\par Wenden wir diesen Satz auf die Formeln 27—32 an, so finden wir unmittelbar
\[
\begin{array}{rl}
33)\quad&Bz-Cy=\cos X\Lambda\mathbin{.}r\BareRoot{mp}\\[4pt]
34)\quad&Cx-Az=\cos Y\Lambda\mathbin{.}r\BareRoot{mp}\\[4pt]
35)\quad&Ay-Bx=\cos Z\Lambda\mathbin{.}r\BareRoot{mp}.
\end{array}
\]
\par Im Perihel gilt \(\mathrm{d}r=0\), \(r-p=-ae(1-e)\), \(x=r\cos XP=a(1-e)\cos XP\), \(y=a(1-e)\cos YP\), \(z=a(1-e)\cos ZP\). Die Formeln 21, 22, 23 lauten daher
\[
\begin{array}{rl}
36)\quad&CG-BH=mpe\cos XP\\[4pt]
37)\quad&AH-CF=mpe\cos YP\\[4pt]
38)\quad&BF-AG=mpe\cos ZP.
\end{array}
\]
\par Verbinden wir sie mit den Gleichungen 30, 31, 32, so liefert unser Lemma
\[
B(BF-AG)-C(AH-CF)=m^{\frac32}p^{\frac32}e\cos X\Pi.
\]
\par Nun ergibt sich für dieselbe Größe \(=F(AA+BB+CC)-A(AF+BG+CH)=mpF\). Somit haben wir
\[
\begin{array}{rl}39)\quad&F=e\BareRoot{mp}\mathbin{.}\cos X\Pi\end{array}
\]
und ebenso
\[
\begin{array}{rl}
40)\quad&G=e\BareRoot{mp}\mathbin{.}\cos Y\Pi\\[4pt]
41)\quad&H=e\BareRoot{mp}\mathbin{.}\cos Z\Pi.
\end{array}
\]
